\documentclass[10pt,letterpaper]{article}
\usepackage{cornell-notes}

\title{Clause 26.08.05: Class Generator Promise Type}
\author{}
\date{}
\setDocTitle{Clause 26.08.05: Class Generator Promise Type}
\setDocSubtitle{C++ 2024}
\setDocOwner{C++ 2024 Cornell Notes}
\setDocDate{\today}
\setCornellCollection{C++ 2024 Cornell Notes}
\setCornellUnitType{Clause}
\setCornellUnitNumber{26.08.05}
\setCornellUnitTitle{Class Generator Promise Type}
\hypersetup{pdftitle={Clause 26.08.05: Class Generator Promise Type},pdfsubject={Cornell notes},pdfauthor={}}

\begin{document}
\maketitle
\begin{CornellOverviewBox}{Overview}
\textbf{Classification:} Standard library - Normative\\[2pt]
\textbf{Learning objective:} Understand the rules, terminology, and semantic consequences associated with Class generator::promise\_type.
\end{CornellOverviewBox}\begin{CornellSummaryBox}{Summary}
{\sffamily\bfseries\color{Accent} Summary}\par\vspace{3pt}Section 26.8.5 focuses on Class generator::promise\_type. Apply the specified interface together with its mandates, effects, result conditions, exception behavior, and complexity constraints. When applying it, preserve the distinction between mandatory requirements, recommendations, permissions, and behavior deliberately left undefined, unspecified, or implementation-defined.
\end{CornellSummaryBox}

\end{document}
